\section{Count Occurrences of an Element in a Sorted Array}
\label{sec:bs-count-occurrences}

\problemstatement{We are given a sorted array $\textit{arr}$ of $n$
elements, possibly with duplicates, and a target $X$. We must count how
many times $X$ occurs in $\textit{arr}$.}

\begin{patternbox}
\emph{``Count occurrences of $X$ in a sorted array''} is immediately
solvable once you recall that all occurrences of $X$ occupy a single
contiguous block (Section~\ref{sec:bs-first-last}), bounded by the lower
bound and upper bound of $X$. The keyword to internalize: any time a
sorted-array problem asks ``how many,'' check first whether the answer is
simply the \emph{width} of a contiguous range you can find with two
boundary searches, rather than something requiring an actual scan.
\end{patternbox}

\subsection*{Understanding the problem carefully}

The lower bound of $X$ is the first index where $X$ could occur; the upper
bound of $X$ is the first index strictly past where $X$ can occur. Every
index in between holds a value equal to $X$ (since the array is sorted and
both boundaries are tight), and no index outside that range does. So the
count of occurrences is simply the number of indices in
$[\textit{lowerBound}(X), \textit{upperBound}(X))$, i.e.
\[
  \textit{count} = \textit{upperBound}(X) - \textit{lowerBound}(X).
\]
This formula is correct even when $X$ is entirely absent: in that case
$\textit{lowerBound}(X) = \textit{upperBound}(X)$ (both point to the same
insertion position), giving a count of $0$ automatically, with no special
case needed.

\begin{ideabox}[Why No Presence Check Is Needed Here]
Unlike Section~\ref{sec:bs-first-last}, where we needed an explicit
presence check before reporting indices (an absent element has no valid
index to report), here the quantity we want is a \emph{count}, and the
subtraction $\textit{upperBound}(X) - \textit{lowerBound}(X)$ naturally
evaluates to $0$ when $X$ is absent, without requiring a separate branch.
\end{ideabox}

\subsection*{Algorithm}

\begin{enumerate}
  \item $\textit{lb} \gets \textit{lowerBound}(\textit{arr}, X)$.
  \item $\textit{ub} \gets \textit{upperBound}(\textit{arr}, X)$.
  \item Return $\textit{ub} - \textit{lb}$.
\end{enumerate}

\begin{algorithm}[H]
\caption{Count Occurrences}
\begin{algorithmic}[1]
\Procedure{CountOccurrences}{$\textit{arr}, X$}
  \State $\textit{lb} \gets \Call{LowerBound}{\textit{arr}, X}$
  \State $\textit{ub} \gets \Call{UpperBound}{\textit{arr}, X}$
  \State \Return $\textit{ub} - \textit{lb}$
\EndProcedure
\end{algorithmic}
\end{algorithm}

\subsection*{Dry run}

\dryrun{Take $\textit{arr} = [2, 4, 4, 4, 4, 8, 10]$, $X = 4$.}

$\textit{lowerBound}(4)$: first index with value $\ge 4$ is index $1$.
$\textit{upperBound}(4)$: first index with value $> 4$ is index $5$
($\textit{arr}[5]=8$). So $\textit{count} = 5 - 1 = 4$, correctly matching
the four occurrences of $4$ at indices $1,2,3,4$.

Now take $X = 6$ on the same array (absent). Both
$\textit{lowerBound}(6)$ and $\textit{upperBound}(6)$ equal $5$ (the first
index with value $\ge 6$, and also the first index with value $>6$, is
index $5$, where $\textit{arr}[5]=8$). So $\textit{count} = 5-5=0$,
correctly reporting absence.

\subsection*{C++ implementation}

\begin{lstlisting}
#include <bits/stdc++.h>
using namespace std;

int lowerBound(vector<int>& arr, int X) {
    int n = arr.size();
    int lo = 0, hi = n - 1, ans = n;
    while (lo <= hi) {
        int mid = lo + (hi - lo) / 2;
        if (arr[mid] >= X) { ans = mid; hi = mid - 1; }
        else lo = mid + 1;
    }
    return ans;
}

int upperBound(vector<int>& arr, int X) {
    int n = arr.size();
    int lo = 0, hi = n - 1, ans = n;
    while (lo <= hi) {
        int mid = lo + (hi - lo) / 2;
        if (arr[mid] > X) { ans = mid; hi = mid - 1; }
        else lo = mid + 1;
    }
    return ans;
}

int countOccurrences(vector<int>& arr, int X) {
    return upperBound(arr, X) - lowerBound(arr, X);
}

int main() {
    vector<int> arr = {2, 4, 4, 4, 4, 8, 10};
    cout << "Occurrences of 4: " << countOccurrences(arr, 4) << endl;
    cout << "Occurrences of 6: " << countOccurrences(arr, 6) << endl;
    return 0;
}
\end{lstlisting}

\begin{complexitybox}
\textbf{Time Complexity.} Two binary searches, each $O(\log n)$, giving
\[
  O(\log n),
\]
a dramatic improvement over the $O(n)$ linear scan that counting
occurrences would otherwise require.
\textbf{Space Complexity.} $O(1)$ auxiliary space.
\end{complexitybox}

\begin{takeawaybox}
``How many elements satisfy property $P$'' on a sorted array, where $P$
itself is monotonic (true on a contiguous range), almost always reduces to
``find the two boundaries of the true range and subtract.'' This is a
strictly more general restatement of the lower/upper-bound subtraction
trick, worth recognizing beyond just counting exact value matches.
\end{takeawaybox}
